\documentclass[11pt]{article}
\usepackage{amsmath,amssymb,amsthm}
\usepackage[hidelinks]{hyperref}
\newtheorem{result}{Result}
\newtheorem{conjecture}{Conjecture}
\newtheorem{theorem}{Theorem}
\newcommand{\card}[1]{\textsuperscript{\texttt{[#1]}}}

\title{Working draft}
\author{}
\date{}

\begin{document}
\maketitle

\begin{abstract}
\end{abstract}

\section{Introduction}
% The question and why it matters: to write when you are ready. This draft can
% stay as it is until a session gives you a result worth keeping.

% This draft is the best current state of the research. Only the lead decides
% what enters it: by ticking a line in a debrief ("Copy Result 1 of note S2 into the
% draft, as a result"), or by editing this file. A result names its claim card, and
% that is what makes the claim established:
%
%   \begin{result}[Short title]
%   The statement, with its key equation.\footnote{Derivation D-S2-1 and computation
%   C-S2-1, each confirmed by a blind review.}\card{CL-S2-1}
%   \end{result}
%
% A claim without a confirming review goes in a conjecture environment instead.
% Derivations behind the results go in appendices, after \appendix.

\section{The period to second order in the amplitude}

% Copied from note S1 by session S2, as ticked in the S1 debrief.
\begin{result}[Second-order coefficient of the period]
For the simple pendulum at amplitude $\theta_0$,
\begin{equation}
  T(\theta_0) = T_0\left(1 + \frac{\theta_0^2}{16} + O(\theta_0^4)\right),
  \qquad T_0 = 2\pi\sqrt{L/g},
  \label{eq:result}
\end{equation}
so $c_2 = 1/16$.\footnote{Derivation D-S1-1 and the independent computation
C-S1-2, each confirmed by a blind review.}\card{CL-S1-1}
\end{result}

\bibliographystyle{plain}
\bibliography{references}
\end{document}
